% GENERATED FILE. Edit canonical lessons, then run npm run generate:books.
% canonical-source-hash: fnv1a64:084997ac3cf69cee
% canonical-lessons: ML-C353-kimvadanti, ML-C353-abhyuham, ML-C353-vyaja, ML-C353-pracarikkuka, ML-C353-pracarippikkuka

\chapter{Rumours}
\label{ch:ml-rumours}
\hlchaptermodality{353}

\begin{chapteropening}
\textbf{By the end of this chapter:} \emph{I can say a rumour, speculation, fake, to circulate and to spread around.}

Complete the last lesson of chapter 353: I can say a rumour, speculation, fake, to circulate and to spread around.
\end{chapteropening}

\section[\texorpdfstring{\ml{കിംവദന്തി} (kiṁvadanti)}{കിംവദന്തി (kiṁvadanti)}]{\ml{കിംവദന്തി} (kiṁvadanti) — a rumour}

\label{lesson:ML-C353-kimvadanti}



\begin{quote}
Before the new one: say the Malayalam for to exaggerate, then the Malayalam for to say sorry.
\end{quote}

\subsection*{You'll want to know: \ml{കിംവദന്തി}}
\textbf{\ml{കിംവദന്തി}} — \emph{kiṁvadanti} — \textquotedblleft{}a rumour\textquotedblright{}.

\begin{quote}
\ml{ആ} \ml{കിംവദന്തി} \ml{സത്യമല്ല} \ml{എന്ന്} \ml{പോലീസ്} \ml{പറഞ്ഞു}.
\end{quote}

\begin{quote}
\emph{ā kiṁvadanti satyamalla ennŭ pōlīsŭ paṟaññu.}
\end{quote}

\begin{quote}
The police said that the rumour is not true.
\end{quote}

A written word for talk that passes from mouth to mouth.

\begin{cousinweb}
\textbf{\ml{കിംവദന്തി}} (\emph{kiṁvadanti}) is borrowed from Sanskrit \emph{kiṁvadantī}, \textquotedblleft{}common talk, a rumour\textquotedblright{} (Monier-Williams), a word made from \emph{kim}, \textquotedblleft{}what\textquotedblright{}, and \emph{vadanti}, \textquotedblleft{}they say\textquotedblright{}: a rumour is what people are saying.
\end{cousinweb}

\begin{grammarlens}[title={What you've built}]
One more word for telling what happened.
\end{grammarlens}

\subsection*{Your turn}
\begin{itemize}
\raggedright
  \item \emph{Say it:} \emph{kiṁvadanti}
  \item \emph{Say it:} \emph{kiṁvadanti}, once more
  \item \emph{Say it:} \emph{ā kiṁvadanti satyamalla ennŭ pōlīsŭ paṟaññu.} — The police said that the rumour is not true.
  \item \emph{Recall:} say the Malayalam for to exaggerate, then the Malayalam for to say sorry, then say \emph{kiṁvadanti} again
\end{itemize}

\begin{tcolorbox}[breakable,colback=teal!4,colframe=teal!35!black,title={Before you move on}]
What does \ml{കിംവദന്തി} mean? (\textquotedblleft{}A rumour\textquotedblright{}.) Say it once more.
\end{tcolorbox}
\section[\texorpdfstring{\ml{അഭ്യൂഹം} (abhyūhaṁ)}{അഭ്യൂഹം (abhyūhaṁ)}]{\ml{അഭ്യൂഹം} (abhyūhaṁ) — speculation; a guess}

\label{lesson:ML-C353-abhyuham}



\begin{quote}
Before the new one: say the Malayalam for to say sorry, then the Malayalam for a rumour.
\end{quote}

\subsection*{You'll want to know: \ml{അഭ്യൂഹം}}
\textbf{\ml{അഭ്യൂഹം}} — \emph{abhyūhaṁ} — \textquotedblleft{}speculation; a guess\textquotedblright{}.

\begin{quote}
\ml{കമ്പനി} \ml{അടയ്ക്കും} \ml{എന്ന്} \ml{അഭ്യൂഹം} \ml{ഉണ്ട്}.
\end{quote}

\begin{quote}
\emph{kampani aṭaykkuṁ ennŭ abhyūhaṁ uṇṭŭ.}
\end{quote}

\begin{quote}
There is speculation that the company will close.
\end{quote}

Plural: \ml{അഭ്യൂഹങ്ങൾ}, \emph{abhyūhaṅṅaḷ}.

In the news the plural is common: \ml{അഭ്യൂഹങ്ങൾ}, speculation.

\begin{cousinweb}
\textbf{\ml{അഭ്യൂഹം}} (\emph{abhyūhaṁ}) is borrowed from Sanskrit \emph{abhyūha}, \textquotedblleft{}reasoning, inference, conjecture\textquotedblright{} (Monier-Williams). In Malayalam news it is a guess that is going round without proof.
\end{cousinweb}

\begin{grammarlens}[title={What you've built}]
One more word for telling what happened.
\end{grammarlens}

\subsection*{Your turn}
\begin{itemize}
\raggedright
  \item \emph{Say it:} \emph{abhyūhaṁ}
  \item \emph{Say it:} \emph{abhyūhaṁ}, once more
  \item \emph{Say it:} \emph{kampani aṭaykkuṁ ennŭ abhyūhaṁ uṇṭŭ.} — There is speculation that the company will close.
  \item \emph{Recall:} say the Malayalam for to say sorry, then the Malayalam for a rumour, then say \emph{abhyūhaṁ} again
\end{itemize}

\begin{tcolorbox}[breakable,colback=teal!4,colframe=teal!35!black,title={Before you move on}]
What does \ml{അഭ്യൂഹം} mean? (\textquotedblleft{}Speculation; a guess\textquotedblright{}.) Say it once more.
\end{tcolorbox}
\section[\texorpdfstring{\ml{വ്യാജ} (vyāja)}{വ്യാജ (vyāja)}]{\ml{വ്യാജ} (vyāja) — fake; forged}

\label{lesson:ML-C353-vyaja}



\begin{quote}
Before the new one: say the Malayalam for a rumour, then the Malayalam for speculation.
\end{quote}

\subsection*{You'll want to know: \ml{വ്യാജ}}
\textbf{\ml{വ്യാജ}} — \emph{vyāja} — \textquotedblleft{}fake; forged\textquotedblright{}.

\begin{quote}
\ml{അവന്റെ} \ml{കയ്യിൽ} \ml{ഒരു} \ml{വ്യാജ} \ml{സർട്ടിഫിക്കറ്റ്} \ml{ഉണ്ടായിരുന്നു}.
\end{quote}

\begin{quote}
\emph{avanṟe kayyil oru vyāja sarṭṭiphikkaṟṟŭ uṇṭāyirunnu.}
\end{quote}

\begin{quote}
He had a fake certificate.
\end{quote}

It stands before the noun, unchanged: \ml{വ്യാജ} \ml{വാർത്ത}, fake news.

\begin{cousinweb}
\textbf{\ml{വ്യാജ}} (\emph{vyāja}) is borrowed from Sanskrit \emph{vyāja}, \textquotedblleft{}deceit, fraud; a disguise, a pretext\textquotedblright{} (Monier-Williams). Malayalam uses the bare stem before a noun, as it does with \ml{പൊതു} and \ml{സ്വകാര്യ}.
\end{cousinweb}

\begin{grammarlens}[title={What you've built}]
One more word for telling what happened.
\end{grammarlens}

\subsection*{Your turn}
\begin{itemize}
\raggedright
  \item \emph{Say it:} \emph{vyāja}
  \item \emph{Say it:} \emph{vyāja}, once more
  \item \emph{Say it:} \emph{avanṟe kayyil oru vyāja sarṭṭiphikkaṟṟŭ uṇṭāyirunnu.} — He had a fake certificate.
  \item \emph{Recall:} say the Malayalam for a rumour, then the Malayalam for speculation, then say \emph{vyāja} again
\end{itemize}

\begin{tcolorbox}[breakable,colback=teal!4,colframe=teal!35!black,title={Before you move on}]
What does \ml{വ്യാജ} mean? (\textquotedblleft{}Fake; forged\textquotedblright{}.) Say it once more.
\end{tcolorbox}
\section[\texorpdfstring{\ml{പ്രചരിക്കുക} (pracarikkuka)}{പ്രചരിക്കുക (pracarikkuka)}]{\ml{പ്രചരിക്കുക} (pracarikkuka) — to circulate; to go round}

\label{lesson:ML-C353-pracarikkuka}



\begin{quote}
Before the new one: say the Malayalam for speculation, then the Malayalam for fake.
\end{quote}

\subsection*{You'll want to know: \ml{പ്രചരിക്കുക}}
\textbf{\ml{പ്രചരിക്കുക}} — \emph{pracarikkuka} — \textquotedblleft{}to circulate; to go round\textquotedblright{}.

\begin{quote}
\ml{ആ} \ml{വാർത്ത} \ml{പെട്ടെന്ന്} \ml{ഗ്രാമം} \ml{മുഴുവൻ} \ml{പ്രചരിച്ചു}.
\end{quote}

\begin{quote}
\emph{ā vārtta peṭṭennŭ grāmaṁ muḻuvan pracariccu.}
\end{quote}

\begin{quote}
The news quickly went round the whole village.
\end{quote}

Past: \ml{പ്രചരിച്ചു}, \emph{pracariccu}.

Said of news and stories: they circulate by themselves.

\begin{cousinweb}
\textbf{\ml{പ്രചരിക്കുക}} (\emph{pracarikkuka}) is built from Sanskrit \emph{pracāra}, \textquotedblleft{}roaming about; coming forth, appearing\textquotedblright{} (Monier-Williams), with the Malayalam verb ending -\ml{ിക്കുക}. What circulates goes about from person to person.
\end{cousinweb}

\begin{grammarlens}[title={What you've built}]
One more word for telling what happened.
\end{grammarlens}

\subsection*{Your turn}
\begin{itemize}
\raggedright
  \item \emph{Say it:} \emph{pracarikkuka}
  \item \emph{Say it:} \emph{pracarikkuka}, once more
  \item \emph{Say it:} \emph{ā vārtta peṭṭennŭ grāmaṁ muḻuvan pracariccu.} — The news quickly went round the whole village.
  \item \emph{Recall:} say the Malayalam for speculation, then the Malayalam for fake, then say \emph{pracarikkuka} again
\end{itemize}

\begin{tcolorbox}[breakable,colback=teal!4,colframe=teal!35!black,title={Before you move on}]
What does \ml{പ്രചരിക്കുക} mean? (\textquotedblleft{}To circulate; to go round\textquotedblright{}.) Say it once more.
\end{tcolorbox}
\section[\texorpdfstring{\ml{പ്രചരിപ്പിക്കുക} (pracarippikkuka)}{പ്രചരിപ്പിക്കുക (pracarippikkuka)}]{\ml{പ്രചരിപ്പിക്കുക} (pracarippikkuka) — to spread around}

\label{lesson:ML-C353-pracarippikkuka}



\begin{quote}
Before the new one: say the Malayalam for fake, then the Malayalam for to circulate.
\end{quote}

\subsection*{You'll want to know: \ml{പ്രചരിപ്പിക്കുക}}
\textbf{\ml{പ്രചരിപ്പിക്കുക}} — \emph{pracarippikkuka} — \textquotedblleft{}to spread around\textquotedblright{}.

\begin{quote}
\ml{ആരോ} \ml{ഒരു} \ml{വ്യാജ} \ml{വാർത്ത} \ml{പ്രചരിപ്പിച്ചു}.
\end{quote}

\begin{quote}
\emph{ārō oru vyāja vārtta pracarippiccu.}
\end{quote}

\begin{quote}
Someone spread a fake news story.
\end{quote}

Past: \ml{പ്രചരിപ്പിച്ചു}, \emph{pracarippiccu}.

The -\ml{ിപ്പിക്കുക} form makes someone else do it: \ml{പ്രചരിക്കുക}, to go round; \ml{പ്രചരിപ്പിക്കുക}, to send round.

\begin{cousinweb}
\textbf{\ml{പ്രചരിപ്പിക്കുക}} (\emph{pracarippikkuka}) is the causative of \ml{പ്രചരിക്കുക} (\emph{pracarikkuka}), \textquotedblleft{}to circulate\textquotedblright{}, which is built from Sanskrit \emph{pracāra}, \textquotedblleft{}going about, appearing\textquotedblright{} (Monier-Williams). The -\ml{പ്പിക്കുക} ending turns going round into sending round.
\end{cousinweb}

\begin{grammarlens}[title={What you've built}]
One more word for telling what happened.
\end{grammarlens}

\subsection*{Your turn}
\begin{itemize}
\raggedright
  \item \emph{Say it:} \emph{pracarippikkuka}
  \item \emph{Say it:} \emph{pracarippikkuka}, once more
  \item \emph{Say it:} \emph{ārō oru vyāja vārtta pracarippiccu.} — Someone spread a fake news story.
  \item \emph{Recall:} say the Malayalam for fake, then the Malayalam for to circulate, then say \emph{pracarippikkuka} again
\end{itemize}

\begin{tcolorbox}[breakable,colback=teal!4,colframe=teal!35!black,title={Before you move on}]
What does \ml{പ്രചരിപ്പിക്കുക} mean? (\textquotedblleft{}To spread around\textquotedblright{}.) Say it once more.
\end{tcolorbox}
